\documentclass[pdflatex,sn-mathphys-num]{sn-jnl}

\usepackage{graphicx}%
\usepackage{multirow}%
\usepackage{amsmath,amssymb,amsfonts}%
\usepackage{amsthm}%
\usepackage{mathrsfs}%
\usepackage[title]{appendix}%
\usepackage{xcolor}%
\usepackage{textcomp}%
\usepackage{manyfoot}%
\usepackage{booktabs}%
\usepackage{algorithm}%
\usepackage{algorithmicx}%
\usepackage{algpseudocode}%
\usepackage{listings}%

\usepackage{caption}
\usepackage{svg}
\usepackage{hyperref}
\newcommand{\figpanel}[2]{\hyperref[#1]{Fig.~\ref*{#1}#2}}
\newcommand{\tabpanel}[1]{\hyperref[#1]{Tab.~\ref*{#1}}}

\raggedbottom

\begin{document}

%% Header
\title{Estimating Neural Timescales}

\author*[1]{\fnm{Ryan} \sur{Hammonds} \email{rphammonds@ucsd.edu}}
\author[2]{\fnm{Dillan} \sur{Cellier}}
\author[3]{\fnm{Blanca} \sur{Martin-Burgos}}
\author[3]{\fnm{Trevor} \sur{McPherson}}
\author[4]{\fnm{Richard} \sur{Gao}}
\author*[1, 2, 3, 5]{\fnm{Bradley} \sur{Voytek} \email{bradley.voytek@gmail.com}}

\affil[1]{\orgdiv{Halıcıoğlu Data Science Institute}}
\affil[2]{\orgdiv{Department of Cognitive Science}}
\affil[3]{\orgdiv{Neurosciences Graduate Program}}
\affil[4]{\orgdiv{Kavli Institute for Brain and Mind}, \orgname{University of California, San Diego}, \city{La Jolla}, \postcode{92093}, \state{CA}, \country{USA}}

\affil[5]{\orgdiv{Bernstein Center for Computational Neuroscience}}
\affil[6]{\orgdiv{Machine Learning in Science}, \orgname{University of Tübingen \& Tübingen AI Center}, \state{Tübingen}, \country{Germany}}

\abstract{
In our collective attempt to understand the ``neural code'', we have developed many conceptual and mathematical tools for deciphering neuronal activity. In sensory cortices, all forms of neural activity, from spiking to EEG, can reliably predict future neural activity up to a few hundred milliseconds later, after which neural activity is more independent. In contrast, activity in association cortex predicts subsequent activity up to several seconds later. This temporal dependency is often captured under the conceptual umbrella of ``neural timescales'', a concept that has received a lot of recent attention because signals with longer timescales are argued to have a longer ``memory'', making timescales attractive for neuroscience applications. While intuitive, a timescale is never measured directly. It is estimated as a parameter of a model, usually an exponential fit to the autocorrelation or a Lorentzian fit to the power spectrum, and it inherits every assumption of that model. Here we ask when a single timescale is recoverable at all. Using simulations with known ground truth, we show that when the data really are first-order autoregressive all common estimators agree to within a few percent, and that when they are not, the same estimators disagree by more than an order of magnitude while each remaining perfectly reproducible. We then show that higher-order autoregressive, AR($p$), models can be read directly as a small set of decay times and rhythms, rather than as an uninterpretable list of coefficients, and that this recovers constructed timescales that the coefficients themselves cannot express. We derive how much data a timescale estimate needs, finding that its precision depends on the recording only through the ratio of recording length to timescale. Finally, we apply one identical pipeline to eight recordings spanning fMRI, MEG, EEG, ECoG, LFP and spiking activity in four species. In every recording a first-order model was inadequate, the estimators disagreed by factors of 1.3 to 179, and between one and eleven autoregressive modes were required to describe the spectrum. We provide recommendations for reporting neural timescales in a way that makes these assumptions visible.
}

\maketitle

\section{Timescales}\label{sec:timescales}

Many real world processes are correlated at varying scales. For example, when trying to predict tomorrow's high temperature, a reasonable starting point is to use the high temperature from today. This also works to predict the high temperature the day after, and the day after that, and so on. One common approach for capturing these temporal relationships is through the use of an autocorrelation, which quantifies how correlated a signal is with itself at every temporal lag. The autocorrelation function (ACF) describes the resulting curve, which can be visualized as the correlation at each time lag (\figpanel{fig:ar1}{}).

At very short lags the correlation is high, and the farther out in time you get, the lower the correlations become. The rate at which the ACF decays from correlated to uncorrelated is usually described as the ``timescale'' of the process. For purely random, Gaussian distributed, white noise processes the timescale is very fast, dropping from perfectly correlated at lag-zero to uncorrelated at all other lags. Timescales are often quantified as the exponential decay rate of the ACF, where numerous real-world processes, ranging from the atomic scale with radioactive decay \cite{semkow_oscillations_2009} to the astronomical with stellar flares \cite{anfinogentov_decaying_2013}, show temporal dependence. In the context of neuroscience, timescales have been observed across recording modalities, from single unit recordings \cite{murray_hierarchy_2014} and neuronal population local field potentials (LFP) \cite{haider_millisecond_2016} to human electroencephalography (EEG) \cite{meisel_decline_2017}, magnetoencephalography (MEG) \cite{zhigalov_relationship_2015}, and functional magnetic resonance imaging (fMRI) \cite{watanabe_atypical_2019, raut_organization_2020}. Regardless of the domain or modality, the mathematical basis of timescales is the same, while interpretations and ground truth mechanisms vary based on the signal's context.

Returning to the temperature example, if we were to record temperature at much higher fidelity we would see that temperature fluctuates periodically on a day/night basis, as well as more slowly on a seasonal basis. These periodic oscillations introduce oscillations in the ACF, because an oscillation is strongly anticorrelated with itself at every half-cycle lag and strongly correlated at every full-cycle lag. Just like with temperature, neural data also oscillates periodically, though often manifesting as more infrequent bursts. Neural oscillations are embedded within a non-oscillatory, aperiodic signal, which is thought to reflect asynchronous neuronal drive and to capture the timescale of the underlying neural activity \cite{gao_neuronal_2020}. This mixing of aperiodic and periodic features means that the ACF also blends aperiodic timescales with neural oscillations. As we explore in more detail below, this mixing can result in significant biases and errors in estimates of the timescale.

To overcome this mixing of signal features, we previously leveraged the mathematical relationship between the ACF and the neural power spectral density (PSD) to show that aperiodic neuronal timescales can be more reliably estimated from the PSD, independent of the influence of oscillations. To do this we used our power spectral parameterization approach for separating periodic and aperiodic processes, where oscillations manifest as narrowband ``bumps'' atop the aperiodic spectral component from which the neural timescale is derived \cite{donoghue_parameterizing_2020}. Using this method, we showed that neural timescales in human participants are dynamic depending on task state, dilating or expanding during working memory \cite{gao_neuronal_2020}.

While the spectral approach helps overcome some of the biases introduced by oscillations, short-time spectral estimates are noisy and limited by their small sampling duration, especially within the context of a behavioral or cognitive task. This finite sampling bias also significantly impacts ACF-based timescale estimates. These three factors -- the presence of oscillations, the fact that timescales can themselves be dynamic as a function of behavioral state, and the finite sampling bias -- all have significant influences on any given timescale estimate. That is, most common methods for quantifying neuronal timescales suffer from some kind of bias.

In this paper we ask a question that comes before the choice of estimator: when is a single neural timescale recoverable at all, and what does an estimate mean when it is not? We illustrate how the common methods are impacted by oscillations, finite samples, dynamic shifts in timescales, and model mismatch. We show this in both simulated and real data, and provide recommendations for addressing these issues. Finally, at the risk of introducing yet another method with implicit and overlooked estimation biases, we develop autoregressive (AR) models as a framework that makes each of these assumptions explicit and testable, and we show how to read such a model as a small set of decay times and rhythms rather than as a list of coefficients.

\section{Autoregressive Models}\label{sec:ar}

The simplest AR model is the first-order autoregressive, AR(1), model: $T_{\text{tomorrow}} = \varphi T_{\text{today}} + \epsilon$. Assuming tomorrow's temperature will be very close to today's, the coefficient $\varphi$ will approach one. As $\varphi$ approaches zero, today's temperature is increasingly less indicative of tomorrow's temperature and trends towards a Gaussian distributed random white noise process dominated by $\epsilon$.

The AR(1) model may be applied to varying temporal scales. If $T_0$ is today's temperature and $T_k$ is the temperature in $k$ days, we can update the temperature model to $T_k = \varphi^k T_0 + \epsilon$. This AR(1) form is a discrete-time analog of the stochastic motion of particles under Brownian motion, replacing temperature with velocity \cite{uhlenbeck_theory_1930}. In these lagged cases the AR(1) coefficient depends on the discrete time step $k$, and is equivalent to the value of the ACF at the corresponding lag, given the sampling rate $f_s$. The ACF of an AR(1) process is an exponentially decaying function and may be expressed in terms of $\varphi$ (eq. \ref{eq:acf}). The autocorrelation at lag one is equal to $\varphi$. The AR(1) model can simulate a signal $x$ given a white noise input (eq. \ref{eq:ar}).

\begin{align}
x_t &= \varphi x_{t-1} + \epsilon \label{eq:ar} \\
A(k) &= e^{\frac{-k}{\tau f_s}} = \varphi^k \label{eq:acf} \\
\tau &= \frac{-1}{f_s \log{\varphi}} \label{eq:tau}
\end{align}

Timescales, $\tau$, may therefore be expressed as the AR(1) coefficient (eq. \ref{eq:tau}), and are bounded by the scale of the observation, defined by the sampling rate $f_s$ and the length of the signal $n$. Sampling rate lower bounds observable timescales, as very short timescales relative to $f_s$ will be undetectable. For example, if a timescale related to temperature fluctuations is on the order of days, but temperature is sampled once a week, the daily timescale will be impossible to detect. Inversely, the length of the signal upper bounds detectable timescales. If a timescale is on the order of years, but only one week of data is collected, the yearly timescale will be undetectable.

\section{Neural Timescales}\label{sec:neural}

Neural timescales characterize the length of time over which an input is integrated and maintained by neuronal populations. That is, if an input causes a change in neural activity that then immediately reverts back to random, its timescale will be very short, or fast: you cannot predict future values from the current state. However, if an input causes a lasting, predictable structure in activity, its timescale will be longer, or slower. This is interpreted as that activity having a longer ``memory'' of the input, which is what makes timescales particularly attractive from a computational neuroscience perspective. However, this is also what makes accurate measurement so critical, especially given the rapid emergence of studies that examine the functional role of timescales in cognition \cite{cavanagh_diversity_2020, gao_neuronal_2020} and their alteration in disease states \cite{watanabe_atypical_2019, wengler_distinct_2020}.

Related to the potential functional role of neural timescales is the fact that they are not randomly distributed across the brain. Rather, they are spatially organized in a cortical hierarchy \cite{hasson_hierarchy_2008, kiebel_hierarchy_2008, murray_hierarchy_2014, gao_neuronal_2020, misic_brain_2010, raut_organization_2020}. The lowest levels of the hierarchy are associated with short timescales, on the order of milliseconds, in visual \cite{buracas_efficient_1998} and auditory \cite{yang_millisecond-scale_2008} sensory cortices. At higher levels, longer timescales on the order of seconds are associated with decision making \cite{brody_timing_2003} and working memory \cite{wang_probabilistic_2002} in prefrontal and parietal cortices. This hierarchical distribution has also been observed across species, suggesting evolutionary preservation \cite{buzsaki_scaling_2013}.

\begin{figure}[htbp]
  \centering
  \includesvg[width=0.95\linewidth]{fig_0.svg}
  \caption{\textbf{Three views of the same model.} Timescales, represented as an AR(1) model and process (left), autocorrelation (middle), and power spectral density (right), and the corresponding AR(1) forms (bottom). A key connection is between the AR(1) coefficient and timescales, where $\varphi$ is equal to the autocorrelation at lag one, $A(1)$.}
  \label{fig:ar1}
\end{figure}

Despite the increasing popularity of studying neural timescales, there are many caveats and limitations that significantly impact the accuracy of fitting mathematical models to data which, in turn, impacts the accuracy of timescale estimates and downstream interpretations. This work evaluates and compares timescale estimation approaches via parameterization of the ACF and PSD, with respect to the corresponding AR(1) form (\figpanel{fig:ar1}{}). The ACF and PSD are connected through the Fourier transform, described by the Wiener-Khinchin theorem, and the AR(1) model has a form in each. The $\varphi$ parameter and the PSD of an AR(1) process have a long history \cite{yule_1927, gilman_power_1963} to build upon. The complexity of the model increases as the order, or number of past values $p$, increases, e.g. an AR($p$) model. The extensive history of theoretical analysis, flexible complexity, and connection between $\varphi$ and $\tau$ make AR models a natural target when studying neural timescales.

Because ground truth timescales are unknown in neurobiology, we cannot know which method works under which conditions or why without simulation. Signals are simulated using a white noise input to an AR(1) system, or equivalently a discrete Ornstein-Uhlenbeck process \cite{uhlenbeck_theory_1930} or a branching process \cite{zierenberg_description_2020}. A binary spiking process may be simulated by convolving an exponentially decaying kernel with a Poisson point process \cite{meyer_inferring_2019}, sampling to yield a binary process, and binning \cite{murray_hierarchy_2014}. Regardless of simulation method, the core assumption is that the signal produced will be an AR(1) realization, bounded by $0 < \varphi < 1$. These bounds narrow as the sampling rate or number of samples decreases.

\section{Estimation Bias}\label{sec:bias}

Simulations allow the study of timescale estimation bias when analytical solutions do not exist. Timescale estimation is evaluated when the AR(1) assumption holds and when it fails. Specifically, estimation bias is evaluated under: 1) the presence of neural oscillations \cite{berger_uber_1929, buzsaki_neuronal_2004}, 2) finite duration or a decreasing number of samples in the signal \cite{broersen_finite_2000, schwarz_estimating_1978, zeraati_flexible_2022}, 3) dynamic, non-constant timescales associated with behavioral or task states, or neural regions \cite{gao_neuronal_2020, spitmaan_multiple_2020}, and 4) model mismatch, when a process is better described by a higher-order AR($p$) model \cite{broersen_finite_2000}. General recommendations for accounting for these sources of bias are provided.

\begin{figure}[htbp]
  \centering
  \includesvg[width=0.95\linewidth]{fig_1.svg}
  \caption{\textbf{Four ways a timescale estimate goes wrong.} Timescale estimation bias in autocorrelation (top) and power spectral density (bottom). \textbf{a)} Oscillations correspond to a cosine function in the ACF or a delta function in the PSD. Robust regression of the PSD is an optimal solution. \textbf{b)} As signals have fewer samples, the ACF and PSD deviate randomly from the expected form. Longer signals or more observations narrow the confidence interval around the estimated timescale. \textbf{c)} A non-constant timescale violates the assumption of stationarity and accurate timescale estimation depends on an appropriate windowing scheme. \textbf{d)} If observed data deviates from the expected AR(1) form, there will be model mismatch. Robust regression sometimes offers a solution, otherwise a higher order AR($p$) model is suggested.}
  \label{fig:bias}
\end{figure}

\subsection{Oscillations}

In the context of neural data, timescales are typically estimated from the autocorrelation of LFPs or binned spike counts \cite{murray_hierarchy_2014}. However both LFPs and single-unit spiking often exhibit rhythmic fluctuations. While an interesting signal in their own right, oscillations can significantly impact the accuracy of timescale estimation. Oscillatory activity manifests as a cosine in the ACF \cite{brunel_fast_1999} or as a narrow peak in the PSD \cite{he_scale-free_2014, buzsaki_scaling_2013}. Adding a cosine can account for oscillatory components in an ACF model (\figpanel{fig:bias}{a}), but optimizing the cosine frequency is difficult and non-convex without strong priors. A poorly optimized cosine, or the lack of one, will bias timescale estimates.

Recent methods allow parameterization of oscillatory activity from the PSD \cite{donoghue_parameterizing_2020}, modeled as Gaussian peaks. Either robust regression or explicit oscillatory modeling provides an alternative to conventional ACF fitting. In the frequency domain oscillations are limited to narrow frequency bands, so power outside of the oscillatory bandwidth is minimally affected and the timescale can be recovered. A mixture of oscillations and stochastic processes also arises from AR(2) models \cite{stralrowski_charts_1970}, which produce a damped cosine in the ACF. This distinction, between a rhythm added to a signal and a rhythm that is part of the signal's dynamics, turns out to matter a great deal, and we return to it below.

Oscillations may also be removed from a signal. A narrow-band filter will, in addition to oscillatory power, greedily remove aperiodic power in the narrow band. Singular spectrum analysis is another option \cite{hassani_singular_2021}, a method based on eigendecomposition that depends on combining principal components associated with the stochastic process, or removing oscillatory components. These methods are optimal when oscillations are stationary across the signal. Bursting oscillations are similar to the dynamic, non-stationary biases below (\figpanel{fig:bias}{c}), and depend on an accurate burst detection algorithm, such as applying thresholds to the analytic amplitude from the Hilbert transform \cite{feingold_bursts_2015}.

\subsection{Finite Sampling}

Stochastic processes suffer from a statistical bias related to finite duration recorded signals (\figpanel{fig:bias}{b}). Short duration time series increase estimation bias, limit the maximum detectable timescale, and decrease confidence. As the duration of a time series decreases and as $\varphi \to 1$, \textit{a)} timescale estimates are increasingly biased and \textit{b)} confidence intervals around the estimate widen. Moving towards unbiased estimation is necessary for studying the fast temporal dynamics of task designs. Dynamic timescales have been observed in working memory \cite{wasmuht_intrinsic_2018, gao_neuronal_2020}, decision making \cite{ganupuru_flexibility_2019} and auditory processing \cite{chien_constructing_2020} tasks.

In the context of event-related task designs or neuronal spiking activity, bias is often mitigated by averaging the ACFs of trials or binned spike counts \cite{murray_hierarchy_2014}. When there are multiple neurons, channels, or repeat trials, it is recommended to take the median or mean PSD or ACF across recordings prior to model fitting. Averaging across neuronal populations, regions, and time has been shown to reduce estimation bias attributed to short duration recordings \cite{spitzner_mr_2021}. However, averaging limits the ability to examine between-neuron, region, or trial variability, and it assumes that the signals being averaged share a timescale.

Computing confidence intervals around an estimated $\varphi$ has a long history \cite{mann_statistical_1943}. The least-squares estimate of the variance of the AR(1) coefficient is:

\begin{align}
\mathbb{E}[ \hat{\varphi}^2 ] &= \frac{1 - \varphi^2}{n-1} \label{eq:var}
\end{align}

This measure is increasingly biased as $\varphi \to 1$. Confidence intervals may instead be computed as a weighted average between the least-squares and Cauchy estimators \cite{gallagher_small_2008}. These intervals are useful for determining unbiased upper and lower bounds on an estimated $\varphi$ or timescale (\figpanel{fig:ci}{}), and for determining confidence in short window estimates, which are required when neural signals are non-stationary or when timescales are non-constant.

\begin{figure}[htbp]
  \centering
  \includesvg[width=0.95\linewidth]{fig_2.svg}
  \caption{\textbf{Confidence narrows slowly with data.} Finite bias and confidence intervals (CIs) around estimated AR(1) coefficients. The CIs narrow around the true $\varphi$ as the number of samples increases from 100 (left) to 1000 (right). Estimates and CIs (gray) were averaged over 1000 simulations.}
  \label{fig:ci}
\end{figure}

\subsection{Dynamic Timescales}

Dynamic timescales refer to a signal whose $\varphi$ or $\tau$ is non-constant across time, referred to as a violation of stationarity. This occurs when timescales are associated with behavioral or task states, or with neural regions \cite{gao_neuronal_2020, spitmaan_multiple_2020}. If a signal spans across these states or regions, a single timescale will combine distinct processes and ignore relevant temporal dynamics (\figpanel{fig:bias}{c}). Timescales should be estimated from signals that are relatively homogeneous or stationary, with respect to \textit{a)} anatomical or functional boundaries or \textit{b)} behavioral or task states. This requires selecting an optimal window size and accurately clustering windows, or determining a priori which windows to combine. However, windowing increases finite sampling bias and each estimate becomes less confident. This leads to a min-max problem: minimize violations of stationarity while simultaneously maximizing window size.

\subsection{Model Mismatch}

Model mismatch occurs when the AR(1) model is simply the wrong model, after other biases have been carefully handled (\figpanel{fig:bias}{d}). This is a generalized extension of oscillatory bias, arising when there is temporal structure that does not neatly fall into the AR(1) form. When this occurs, the use of a higher order AR($p$) model offers a solution (eq. \ref{eq:arp}). Increasing the AR order flexibly models the PSD and ACF, allowing for forms that deviate from the expected. AR($p$) models use multiple past values to predict the current value, and also have forms in the ACF and PSD (eq. \ref{eq:arp_psd}).

\begin{align}
    x_t &= \sum_{k=1}^p \varphi_k x_{t-k} + \epsilon \label{eq:arp}\\
    P(f) &= \frac{\sigma^{2}}{|1 - \sum\limits_{k=1}^p \varphi_{k} e\strut^{-j 2 \pi f k / f_s}|^{2}} \label{eq:arp_psd}
\end{align}

Increasing the order comes at two costs. The first is statistical: a more flexible model can absorb other sources of variance, such as finite sampling noise, so the order should be kept as small as the data support. Bayesian information criteria (BIC) may be used to evaluate a set of models with varying $p$ and determine an optimal, minimized order \cite{schwarz_estimating_1978, vidaurre_survey_2013}. The second cost is interpretive, and it is the one that has limited the use of AR($p$) models in this literature. A timescale is a single number, and an AR($p$) fit returns a list of $p$ coefficients that cannot be read as timescales. For $p > 1$, $\varphi_1$ is not an autocorrelation and need not even lie between zero and one, so the usual conversion in eq. \ref{eq:tau} does not apply. We show below that this obstacle is removable.

\section{Results}\label{sec:results}

\subsection{Simulating Known Timescales}\label{sec:sim}

Because the true timescale of a neural recording is unknown, we first ask a question that can be answered exactly: given data generated by a known process, do the standard estimators return the number we put in? We simulated six processes with closed-form autocorrelations and spectra, five of which are shown in \figpanel{fig:benchmark}{a}, and estimated a timescale from each using six estimators. The estimators are organized by the domain in which they are fit -- the time series itself, the autocorrelation, or the spectrum -- and by model order, AR(1) or AR($p$) with the order chosen by BIC. All six consume the same simulated signal, the same ACF and the same spectrum, so any disagreement between them is a property of the model and not of the preprocessing. Recordings were 100 seconds long at 500 Hz with a constructed timescale of 50 ms, which is long enough that finite sampling bias is negligible.

Before comparing estimates to a truth, we have to say what the truth is, because outside AR(1) it is not given. For an AR(1) process $\tau$ is a parameter of the generating equation, and every reasonable definition -- the exponential decay of the autocorrelation, the Lorentzian knee of the spectrum, the area under the autocorrelation -- returns the same number. That coincidence is exactly what fails for the other processes, and it fails differently in each case. The AR(2) oscillator has no relaxation time at all, only an envelope decay. The AR(8) has several timescales, so ``the'' timescale is whichever one a stated rule picks out. The power law has none. And an added sinusoid has one only for its aperiodic part.

We therefore take the target to be whatever the constructed process gives under the \emph{same} rule the estimator is asked to apply, computed from the true poles rather than fitted ones: the relaxation time for AR(1), the envelope decay for the AR(2) oscillator, and for the AR(8) the most important non-rhythmic mode (Methods \ref{sec:methods-est}). This is the only comparison that is not rigged, and it is worth being explicit that it is a choice. A reader who prefers a different definition of ``the'' timescale of a multi-timescale process would get different numbers from the same fits, which is itself part of the argument of this paper.

When the generating process really is AR(1), the choice of estimator does not matter (\figpanel{fig:benchmark}{c}, first column). All six recovered the 50 ms timescale to within 2\%, order selection returned $p=1$ in all three domains, and all six left under 1\% of the log spectrum unexplained. This is the situation the field implicitly assumes, and in it the assumption is harmless.

None of the other five processes behave this way. Adding a 10 Hz sinusoid carrying as much variance as the aperiodic component, a situation that is unremarkable in real EEG or LFP, spread the estimates over a factor of nine: the autocorrelation fit returned 15 ms and the direct time-domain fit 84 ms, while both spectral fits landed within 0.5\% of the truth. The order-selected autocorrelation fit did not return an aperiodic timescale at all. It spent its poles on the rhythm and was left with no relaxation mode, which is why that cell of \figpanel{fig:benchmark}{c} is empty. The reason is visible in the spectrum (\figpanel{fig:benchmark}{b}): a rhythm occupies a handful of frequency bins and can be discounted as an outlier, but it contaminates every lag of the autocorrelation and every term of a lagged regression. Robust regression cannot rescue the lag domain, because in the lag domain the rhythm is not an outlier. For a stochastic AR(2) oscillator, the three AR(1) estimators returned 12, 127 and 229 ms for the same 50 ms process, a nineteen-fold spread, while all three AR($p$) fits recovered the envelope decay to within 3\% and the rhythm to within 0.05 Hz.

The constructed AR(8) separates order from family. It is an autoregressive process, so nothing about it is outside what the model can express; it is simply not first order. The three AR($p$) estimators recovered its aperiodic timescale to within 9\%, returning 50 to 54 ms against a target of 50 ms. The three AR(1) estimators spread over a factor of 401, from 78 ms for the autocorrelation fit to 1.07 s for the lagged regression, while the AR(1) spectral fit ran to the upper bound of the search and returned 20 s, which is the fit reporting as clearly as it can that no single relaxation time describes this spectrum. Failing to raise the order is enough to produce a four-hundred-fold error with no other assumption violated.

A long-memory process, for which no characteristic timescale exists, produced estimates from 4 to 86 ms, a twenty-fold range. Every one of these numbers is stable: repeated simulations return the same value to within a few percent. Reproducibility is not the same as recoverability, and a timescale reported without a model-adequacy check cannot distinguish the two.

To give these cases a common reference we use the integrated timescale, the area under the autocorrelation \cite{sokal_monte_1997}. It equals $\tau$ for an AR(1) process, is the variance-weighted average of the component timescales for a sum of AR(1) processes, and diverges when no timescale exists (\figpanel{fig:benchmark}{c}, orange lines). The AR(2) oscillator is a useful warning in the other direction. Strictly it has no relaxation time at all: its poles are complex, its autocorrelation is a damped cosine rather than a decaying exponential, and there is no lag after which it is simply small. What it has instead is an envelope that decays with a 50 ms time constant, and that is the number we compare estimators against. Its integrated timescale is 9 ms, because the anticorrelated troughs of an oscillatory autocorrelation cancel most of the integral. For an oscillatory process, ``how long does activity stay correlated'' and ``how fast does a perturbation decay'' are different questions with different answers.

Model mismatch is measurable, and it tracks these failures. An AR(1) spectral fit left 0.8\% of the log spectrum unexplained for the AR(1) process, but 2.5\% with an added sinusoid, 24\% for the AR(2) oscillator, 55\% for the AR(8) and 26\% for the power law. AR($p$) reduced these to between 0.1\% and 4.0\%. Order selection tells the same story from the other side: BIC returned $p=1$ for AR(1) and for the added sinusoid, $p=8$ for the AR(8), which is the constructed order, and ran to the top of the search range for the power law. Separating model mismatch from estimation noise confirms that this is not a small-sample effect. Fitting each process's exact autocorrelation, which is what any consistent estimator converges to with infinite data, the mismatch collapses to zero at the correct order for the three processes that are autoregressive and plateaus for the added sinusoid, which no finite order removes.

Because every model is scored in both domains, we can also ask whether the domain a model is fit in changes the model you get. It does, and asymmetrically. Across the four processes that are not AR(1), fits to the spectrum described the autocorrelation nearly as well as fits to the autocorrelation did, leaving between 0.4\% and 7\% of it unexplained, with the added sinusoid the sole exception at 90\%. Fits to the autocorrelation described the spectrum badly, leaving 12\% of it unexplained for the AR(2) oscillator, 7\% for the AR(8), and for the added sinusoid and the long-memory process many times the total spectral variance. The asymmetry is easiest to see in a single fit: on one realization of the long-memory process, the AR(5) model chosen in the lag domain has a spectrum that spans eleven orders of magnitude across the analysis band, where the data span two. The lag-domain objective is simply insensitive to what the model does to the spectrum. This is not an artifact of a poorly optimized autocorrelation fit: both domains use the same model, the same reparameterization, the same order selection and the same robustness setting, the two are given the same band by construction (Methods \ref{sec:methods-est}), and the autocorrelation fit is the better of the two at describing the autocorrelation, leaving under 2\% of it unexplained for every process. A model estimated in the frequency domain is a usable description of the signal in either domain; a model estimated in the lag domain is not.

\begin{figure}[htbp]
  \centering
  \includesvg[width=\linewidth]{revision/figures/fig_benchmark.svg}
  \caption{\textbf{When is a single timescale recoverable?} Five simulated processes, each recorded for 100 s at 500 Hz, 200 repetitions. \textbf{a)} 1.5 s of each simulated signal. \textbf{b)} The power spectrum of one simulation (grey) with an AR(1) fit (dashed) and an order-selected AR($p$) fit (solid). \textbf{c)} The timescale reported by each of six estimators, grouped by the domain they are fit in (time series, autocorrelation, spectrum) and by model order. Each dot is one simulation and the black bar is the median. Three reference lines are drawn where they exist: the constructed relaxation time (black dashed); the decay time of the envelope for the AR(2) oscillator (red dashed), which has no relaxation time because it rings rather than decaying monotonically; and the integrated timescale, the area under the autocorrelation (orange dotted), which is the only one of the three that is defined for every process shown. Where a cell is marked n/a the fitted model contained no aperiodic mode at all. When the data are AR(1) the estimators agree; otherwise they disagree by up to a factor of 46 while each remains individually reproducible.}
  \label{fig:benchmark}
\end{figure}

\subsection{Oscillations Added and Oscillations Intrinsic}\label{sec:osc}

The two oscillatory simulations above fail in opposite ways, and the difference is worth isolating because it determines what can be done about it (\figpanel{fig:osc}{}).

A sinusoid added to an aperiodic signal is not an autoregressive process at all. Its spectrum contains a delta function, which no product of autoregressive poles can reproduce at any finite order, and its autocorrelation contains an undamped cosine that persists at every lag. In the autocorrelation domain this is fatal. The AR(1) estimate falls monotonically as the rhythm grows, from 41 ms when the rhythm is weak to 10 ms when it carries four times the aperiodic power, and reaches an eighteen-fold underestimate at the worst combination of frequency and amplitude we tested (\figpanel{fig:osc}{c}). Raising the model order changes the failure rather than removing it: the AR($p$) autocorrelation fit sits at 60 to 72 ms, an overestimate of 20 to 45\% that does not shrink as the rhythm grows, and across the full grid it ranges from half the truth to twelve times it.

That the higher-order lag-domain fit is also biased deserves an explanation, because the mode decomposition is exact and the natural expectation is that it should therefore be more accurate. The bias is structural rather than a failure of the optimizer or an artifact of finite data. Given the \emph{exact} population autocorrelation of this process, the lag-domain answer still ranges from 17 to 57 ms depending on the order fitted, against a truth of 50 ms. The reason is visible in eq. \ref{eq:acfmodes}: an AR($p$) autocorrelation is a sum of one term per mode, but the weights $A_j$ are functions of the poles rather than free parameters. The data here is an external mixture, $w\,e^{-k/\tau} + (1-w)\cos(\omega k)$, whose weight $w$ is independent of the shape of either component, and the model has no parameter for it. Reproducing an undamped cosine requires a pole essentially on the unit circle, which carries a large residue, and the only way to recover the correct balance is to displace the remaining poles; the aperiodic ones absorb the entire error. Robust regression is available in the lag domain and does not help, because an additive rhythm is an outlier in frequency and not in lag: it perturbs a handful of spectral bins but every single lag. The spectrum has the escape hatch the lag domain lacks. Robust regression there discounts the narrow peak outright and fits a clean AR(1) shape to what remains, returning 49 ms at every rhythm amplitude. Exactness of the \emph{representation} does not buy identifiability of the \emph{mixture}, and only the spectral domain can decline to model the part of the mixture that does not belong to the process. Both spectral estimates stayed between 0.97 and 1.03 times the truth at every one of the 36 combinations, and were essentially unchanged as the rhythm grew from one eighth to four times the aperiodic power (Supplementary \figpanel{fig:s_grid}{}).

A stochastic oscillator is a different object. An AR(2) process with a complex pair of poles produces a rhythm that drifts in phase and amplitude, which is what neural rhythms actually look like \cite{stralrowski_charts_1970, feingold_bursts_2015}, and it is exactly representable by an autoregressive model of order two. Here the estimators behave in the opposite way (\figpanel{fig:osc}{f}). The AR($p$) spectral fit recovered the decay time to within 3\% and the AR($p$) autocorrelation fit to within 9\%, across a decade of envelope decays and two octaves of rhythm frequency, and both recovered the rhythm frequency to better than 0.4 Hz. The AR(1) estimates were wrong by factors between 0.016 and 133 that themselves depended on both the true decay and the rhythm frequency, so no calibration could remove them. That both domains now succeed here, and that only the spectral domain succeeded for the added sinusoid, is the cleanest statement of the difference between the two kinds of rhythm. This is the practical content of the distinction: a rhythm that is part of the dynamics can be modeled, and a rhythm that is added to the signal can only be avoided.

\begin{figure}[htbp]
  \centering
  \includesvg[width=\linewidth]{revision/figures/fig_oscillations.svg}
  \caption{\textbf{A rhythm added to a signal and a rhythm that is part of it.} \textbf{a-c)} An AR(1) process with a 10 Hz sinusoid added. The spectrum contains a line that no autoregressive model can reproduce (b), and only the spectral fits keep the aperiodic timescale as the rhythm grows (c). \textbf{d-f)} An AR(2) stochastic oscillator. Its rhythm differs from the sinusoid of (a) in every property except its centre frequency. A sinusoid has one fixed amplitude and one phase that advances forever, so its spectrum is a line of zero width, its autocorrelation is an undamped cosine that never decays, and it carries no information beyond three numbers. An AR(2) oscillator is driven by noise at every sample, so its amplitude and phase both wander: it waxes and wanes, occasionally skips a beat, and its autocorrelation is a cosine multiplied by a decaying envelope. Its spectrum is a peak of finite width, set by how fast that envelope decays, and the number of cycles it holds phase for, $f\tau$, is what separates a rhythm you could see in the raw trace from a broad spectral shoulder. This is what neural rhythms look like: bursty and phase-drifting rather than metronomic. It is exactly an autoregressive process, an AR(2) model reproduces its spectrum (e), and both AR($p$) fits recover the decay time of its envelope while the AR(1) fits do not (f). The quantity plotted in (f) is not the relaxation time $\tau$ of panel (c): a damped oscillator has no relaxation time, because its autocorrelation is a damped cosine rather than a decaying exponential. What it has is an envelope, $e^{-t/\tau_{\mathrm{env}}}$, multiplying the cosine, and $\tau_{\mathrm{env}}$ is the decay plotted on both axes. It is the timescale over which a perturbation of the rhythm dies away, and it is $-1/(f_s\log|\lambda|)$ for the mode's pole $\lambda$, the same formula as for a relaxation mode.}
  \label{fig:osc}
\end{figure}

\subsection{Reading an AR($p$) Model}\label{sec:modes}

If a single timescale is often inadequate, the natural response is to raise the model order, and the obstacle is that an AR($p$) fit returns a list of coefficients rather than a timescale. That obstacle is removable, because the coefficients are not the natural description of the model.

Every AR($p$) model is equivalent to a small set of \emph{modes} (Methods \ref{sec:methods-modes}). Each mode is one way the process rings down after an input: a decay, with a time constant $\tau$, and possibly a rhythm at frequency $f$. The response of the process to a single input is the sum of these modes (\figpanel{fig:modes}{a}), and its spectrum is built from them as well (\figpanel{fig:modes}{b}). Formally, the modes are the poles of the model, a real pole giving a pure decay and a complex-conjugate pair giving a damped oscillation, with
\begin{align}
\tau_k = \frac{-1/f_s}{\log|\lambda_k|}, \qquad f_k = \frac{f_s |\arg \lambda_k|}{2\pi}. \label{eq:mode}
\end{align}
An AR($p$) fit can therefore be reported as a short list of $(\tau, f)$ pairs rather than as $[\varphi_1, \dots, \varphi_p]$.

Both kinds of mode have a $\tau$, and it is worth being explicit about what each one means, because they are not the same quantity. A real pole contributes $\lambda^n$ to the response: a plain exponential decay, and its $\tau$ is the whole of it. A conjugate pair contributes $r^n\cos(n\theta + \phi)$, a cosine inside a decaying envelope, and its $\tau$ is the decay of that envelope -- how fast the ringing dies away, not how long the process stays correlated. The same formula, $\tau = -1/(f_s\log|\lambda|)$, gives both, because in both cases it is the modulus of the pole that sets the decay and only the argument that distinguishes a rhythm from a relaxation.

Read in the lag domain this gives the decomposition its plainest form: grouping the conjugate terms of eq. \ref{eq:acfmodes}, \emph{the autocorrelation of any AR($p$) process is a sum of exponential decays and damped cosines}, one per mode, with a real pole giving the first and a conjugate pair the second. An AR(1) is the special case with a single exponential and nothing else. This is the sense in which raising the order does not abandon the exponential picture that motivates timescale estimation in the first place; it adds terms to it.

The term is borrowed from linear systems theory, where the modes of a linear system are the eigenvalues that govern how it relaxes, and it is used in the same sense here. A mode is not the same thing as a model order. An AR($p$) model has exactly $p$ poles, but a real pole is one mode and a complex-conjugate pair is also one mode, because the two poles of a pair always move together and describe a single ringing decay. The number of modes is therefore between $p/2$ and $p$: the constructed AR(8) model in \figpanel{fig:modes}{} has five modes, two pure decays and three rhythms, because each rhythm costs two of the eight orders. Reporting the modes rather than the order is what makes the model readable: the slow mode of that model is a pure decay with $\tau = 100$ ms and no rhythm, while its 10 Hz mode is a conjugate pair that rings at 10 Hz with an envelope decaying at $\tau = 50$ ms. Both are timescales, and they are different timescales of different kinds.

Calling a mode oscillatory says only that its poles are complex, which is a statement about the algebra and not about whether anyone would recognise a rhythm. The quantity that settles that is the product $f\tau$, the number of cycles the mode completes before its envelope decays by $1/e$. The 10 Hz mode above has $f\tau = 0.5$: it rings down before completing a single period, and is better described as a broad shoulder than as an oscillation. We report $f\tau$ alongside every oscillatory mode for this reason. It is a property of the mode alone. Whether the fitted \emph{model} shows a peak at that frequency is a separate question, because modes multiply in the spectrum, so a mildly resonant factor can leave no visible bump when the other factors fall steeply through the same frequencies.

The two domains combine the same modes in different ways, and the difference is what makes a higher-order model confusing to read. In the spectrum the modes \emph{multiply}: the model spectrum is a product of one factor per mode, each depending only on its own pole (\figpanel{fig:modes}{b}). In the autocorrelation they \emph{add}:
\begin{align}
\gamma(k) = \sum_j A_j \lambda_j^{\,k}, \label{eq:acfmodes}
\end{align}
a sum of one exponential or damped cosine per mode (\figpanel{fig:modes}{c}). Both are exact, and in our implementation both reproduce the model to better than one part in $10^{10}$.

They are not the same decomposition, and it is worth saying why, because the natural guess is wrong. A product in the frequency domain is a \emph{convolution} in the lag domain, not a sum, so a mode's additive term in eq. \ref{eq:acfmodes} is not the inverse transform of its multiplicative factor in the spectrum; the two differ substantially. The product form and the sum form are the factored and partial-fraction representations of the same rational function. What the partial-fraction form costs is separability: each residue $A_j$ depends on \emph{every} pole, so a mode's contribution to the autocorrelation is not a property of that mode alone.

This decomposition is not new, and we want to be clear that we are importing it rather than deriving it. It is the autoregressive case of a general decomposition of linear time series models into latent components \cite{west_time_1997}, and it has been applied to EEG before, where the latent quasi-cyclical components of a fitted autoregression were used to track the frequency structure of seizures \cite{west_evaluation_1999}. What is new here is the use we put it to: reading a fitted AR($p$) as a set of neural timescales, ranking those modes by how much of the spectrum each accounts for, and asking which of them a recording actually supports.

Not all modes matter equally, so we also rank them. Deleting one mode and refitting the overall scale degrades the fitted spectrum by an amount we can compute in closed form, and we use that degradation as the mode's importance (Methods \ref{sec:methods-importance}). Ranking modes this way lets us ask how many are actually needed: for the constructed AR(8) model built from five modes, three of them reproduce the fitted spectrum to $R^2 = 0.996$, and they are exactly the three substantive ones (\figpanel{fig:modes}{d}). We emphasize that the importance ranking is a statement about the fitted spectrum, not a variance decomposition of the signal. Because all of the modes are driven by the same input, the process variance does not split into per-mode shares, and reporting a mode's ``percent variance explained'' would be wrong.

The modes and the coefficients are not interchangeable, and it is worth being explicit about how far apart they are. Refitting 200 independent 60-second realizations of the same constructed process, every one of the five constructed modes was recovered, with median absolute errors in $\tau$ of 6.5\%, 3.2\%, 2.5\%, 3.8\% and 6.8\% running from the slowest mode to the fastest (\figpanel{fig:modes}{e}). The three rhythm frequencies came back to within 0.9\%: 9.99, 40.1 and 179.7 Hz against 10, 40 and 180. The two relaxations, which are the modes a timescale study would actually report, were recovered to 97.8 and 2.04 ms against 100 and 2.

Neither of the two things one might read instead of a mode has a ground truth to be scored against, which is why neither appears in the figure, but both are worth stating. The first coefficient of the same fits, which a reader might be tempted to convert with eq. \ref{eq:tau}, was $4.04 \pm 0.01$ across the 200 realizations; it exceeds one in every fit and therefore has no timescale at all. An AR(1) fit to the same signals did not converge to a timescale either: its estimate ran to the 20 s upper bound of the search in every realization, which is the fit telling us, as clearly as it can, that no single relaxation time describes this spectrum. Extracting modes is not a refinement of reading coefficients; it is the difference between an interpretable answer and no answer.

\begin{figure}[htbp]
  \centering
  \includesvg[width=\linewidth]{revision/figures/fig_modes.svg}
  \caption{\textbf{An AR($p$) model can be read as a few decay times and rhythms.} A constructed model with five modes: two pure decays, at 100 and 2 ms, and three damped rhythms, at 10, 40 and 180 Hz. The modes were chosen to span the range of timescales and rhythm bands reported in neural recordings; the order follows from that choice, since a real pole costs one order and a conjugate pair costs two, giving $2 + 6 = 8$. \textbf{a)} The response to a single input is the sum of the modes, ordered slowest to fastest and drawn at unit amplitude so their shapes are visible. Blue is a relaxation mode, from a real pole; red is an oscillatory mode, from a conjugate pair; black is the summed response. The same two colours carry the same meaning in every panel. Each oscillatory label gives the cycles it completes before its envelope decays; every one here is under a cycle, so none of them is a rhythm in any useful sense. \textbf{b)} In the spectrum the modes multiply. Grey is one realization, black the generating spectrum, and green the order-selected AR($p$) fit to that realization. \textbf{c)} In the autocorrelation the modes add, as exponential decays for the real poles and damped cosines for the pairs; the amber curve is their sum, which lies on the true autocorrelation exactly. In both panels the individual modes are drawn in their own colour, the three dominant ones dashed and opaque and the two negligible ones dotted and faint. \textbf{d)} Ranking modes by how much the fitted spectrum degrades when each is deleted, three of the five reproduce it to $R^2 = 0.996$. \textbf{e)} Median absolute error in each constructed mode's $\tau$ over 200 independent 60-second realizations, with the modes matched to the construction by pole location. Only quantities with a ground truth appear here, which means modes; all five are recovered to within 7\%.}
  \label{fig:modes}
\end{figure}

\subsection{Neural Signals Across Species and Modalities}\label{sec:realdata}

Simulations show what can go wrong. To ask how often it does, we applied one identical pipeline to eight publicly available recordings spanning fMRI, MEG, EEG, ECoG, LFP and spiking activity, in humans, macaque, mouse and rat (Methods \ref{sec:methods-data}). Each was analysed with the same six estimators and the same mode decomposition, with no dataset-specific artifact rejection, filtering, or oscillation removal. \figpanel{fig:realdata}{} shows the four estimators fit in the two domains that share their inputs exactly, the autocorrelation and the spectrum; the two time-domain fits are tabulated with the rest and quoted where they matter. The aim is not to make a new biological claim about any one recording, but to ask how often the single-timescale model is adequate when the same question is put to signals of very different origin.

All four of the biases sketched in \figpanel{fig:bias}{} were present, and in most recordings they were present together.

\emph{Oscillations.} Every field-potential recording contained a narrowband peak, and the order-selected AR($p$) fit placed a rhythmic mode inside the expected band in six of the eight recordings. The two exceptions are the two that have no narrowband rhythm to find: resting-state fMRI, and the spike counts, whose fitted band lies below the 2-8 Hz range where a rhythm was sought. This is worth stating carefully, because a robust fit is in some tension with the question: robustness exists so that a rhythm does not contaminate the aperiodic timescale, which means the fitted model is under instruction to discount exactly the feature one would then be asking it about. In practice the tension is milder than it sounds, because the extra poles a higher order buys are spent on the rhythm whether or not the loss discounts it.

Whether to be robust is therefore a choice about what is being measured, and it is a trade-off rather than a free lunch. On a process built to test it -- an AR(1) background at $\tau = 50$ ms plus a broadband AR(2) rhythm at 10 Hz, which unlike an added sinusoid the model can represent -- the robust spectral fit recovered the aperiodic timescale best over 20 realizations, at 53 ms against a truth of 50, with AR(1) at 56 ms. Dropping robustness let the model describe the rhythm, which it then found in all 20 fits at a median of 9.4 Hz, and cost 35\% accuracy on the aperiodic timescale, which came back at 68 ms. Fitting the oscillation does not buy the timescale for free; it trades one against the other, and a study that wants both should report them from two fits and say so.

\emph{Model mismatch.} An AR(1) model left between 5.1\% and 41\% of the log spectrum unexplained in the field-potential recordings, 24\% in fMRI and 46\% in the spike counts, with a median across all eight of 14\%. An order-selected AR($p$) model reduced this to between 0.2\% and 3.0\% in the field potentials, 21\% in fMRI and 38\% in the spike counts, with a median across all eight of 1.3\%. Between one and eleven modes were needed to reproduce the fitted spectrum, and the selected order reached the top of the search range in three of the eight recordings. Resting-state fMRI needed only a single mode, and fails the adequacy check below on goodness of fit alone rather than on the number of timescales.

The spike counts are the hardest case in the collection, and they illustrate a point about the analysis band that applies generally. A spike train contributes a flat Poisson floor on top of whatever the membrane dynamics contribute, and the AR spectral model, $\sigma^2/|A|^2$, has no additive term for a floor. Fitted over the full 0.05-8 Hz band the AR(1) model is therefore fitting mostly shot noise: above about 1.5 Hz the binned-count spectrum is within a factor of 1.5 of its own floor, and the least-squares optimum for a single-corner model is essentially flat, leaving 95\% of the log spectrum unexplained. Restricting the fit to 0.05-2 Hz, where the process is still above its shot noise, is what a careful analyst would do and is what we do here; it improves the AR(1) fit tenfold and brings the two spectral estimates from disagreeing by a factor of eleven to a factor of two. Even then the spike counts remain the worst-fit recording in the set, which is a caution for the large literature that estimates timescales from binned spikes: the band over which a spike-count model is fitted is a modelling choice that changes the answer, and it should be reported.

The autocorrelations tell the same story (\figpanel{fig:realdata}{b}): in six of the eight recordings the empirical autocorrelation oscillates rather than decaying monotonically, and only the AR($p$) model follows it. The exceptions are fMRI, whose autocorrelation is over within three samples, and the spike counts, whose fitted band is too narrow to contain a rhythm.

\emph{The consequence.} Because no recording satisfied the AR(1) assumption, the estimators disagreed (\figpanel{fig:realdata}{c}). Among the four estimators plotted, the ratio of the largest to the smallest ranged from 1.3 to 7.3 with a median of 4.7; across all six, including the two time-domain fits, it ranged from 1.3 to 179 with a median of 7.6. The pattern matched the simulations closely. The largest disagreements were in the recordings with the strongest rhythms: motor ECoG at 7.3, posterior alpha at 7.1 and mouse LFP at 6.7, in each case with the AR(1) fits reporting a fast decay and the AR($p$) fits a slower relaxation mode once the rhythm is described rather than absorbed. The single most extreme case involves an estimator the figure omits: in sensorimotor MEG the time-domain fit returned 2.0 s against 11 ms from the autocorrelation and 23 ms from the spectrum, a lagged regression reading a strong sustained beta rhythm as persistence, and it is precisely because a lagged regression fails this way that we do not plot it beside the two domains that share their inputs. The smallest disagreements, under a factor of 1.5, were in resting-state fMRI, where the four estimators returned between 2.3 and 3.0 s, and in macaque ECoG. That fMRI is the most internally consistent recording in the set is worth stating plainly: its spectrum is the one with the least structure to disagree about, not the one estimated best.

Reported on a linear time axis, the estimates also reproduce the scale separation that motivates the hierarchy literature (\figpanel{fig:realdata}{c}). Aperiodic timescales in the field-potential recordings were tens of milliseconds, between 11 and 52 ms by the AR(1) spectral fit, the binned spike counts gave 232 ms, and precuneus BOLD gave 1.0 to 6.6 s across traces. This three-order-of-magnitude gap between hemodynamic and electrophysiological estimates is expected, and it is a useful reminder that comparisons of timescales across modalities are comparisons of different models fit to differently filtered data \cite{raut_organization_2020}.

\begin{figure}[htbp]
  \centering
  \includesvg[width=\linewidth]{revision/figures/fig_realdata.svg}
  \caption{\textbf{The single-timescale model was inadequate in every recording tested.} Eight recordings across six modalities and four species, each analysed with the identical pipeline. \textbf{a)} Power spectrum, median across the traces of that recording (black, solid), with the median AR(1) spectral fit (orange, dashed) and the median order-selected AR($p$) spectral fit (blue, dotted). Rows are labelled by species and modality. Panel c pools every trace, so the curves above it do too. fMRI is estimated by multitaper and the rest by Welch (Methods \ref{sec:methods-est}). \textbf{b)} Autocorrelation over the analysis band, again the median across traces (black, solid), with the AR(1) lag-domain fit (pink, dashed) and the AR($p$) lag-domain fit (green, dotted); the lag axis is cropped to the decay, which differs by three orders of magnitude across rows. Each of the four estimators keeps one colour throughout the figure, so a box in (c) is the colour of the curve it came from. \textbf{c)} The timescale each of those four estimators returns, one box per estimator over all traces of that recording: box is the interquartile range, line the median, whiskers the range. The two time-domain estimators are computed and tabulated but not plotted, so that every curve above has a box below it. For AR($p$) this is the most important non-rhythmic mode, that is, the mode carrying most of the fitted spectrum among those that ring for less than one cycle (Methods \ref{sec:methods-est}); it is not simply the slowest mode, nor the dominant mode overall, which may be a rhythm. The axis is logarithmic because within a single recording the estimators can differ by two orders of magnitude.}
  \label{fig:realdata}
\end{figure}

\subsection{How Much Data Does a Timescale Need?}\label{sec:windows}

The remaining two biases, finite sampling and non-stationarity, are entangled: the standard response to a changing timescale is to shorten the analysis window, and shortening the window is exactly what makes an estimate imprecise. We can now say how the trade-off resolves.

In simulation, the answer for an AR(1) process follows from eq. \ref{eq:finite} and is simple enough to use as a rule of thumb (\figpanel{fig:windows}{b}). Both the bias and the precision of a timescale estimate depend on the recording only through the ratio of duration to timescale, so a bias below 10\% needs about $20\tau$ of data and a standard error below 10\% needs about $200\tau$ (\figpanel{fig:windows}{a,c}). For a 100 ms cortical timescale that is 2 and 20 seconds respectively, per estimate. Precision improves only as the square root of duration, so buying another factor of two costs four times the data.

When there is no timescale to recover, shortening the window does something worse than adding noise. Applied to the long-memory process, the autocorrelation fit returns an estimate proportional to the length of the recording, tracking $\hat\tau \propto T$ over three decades, while the same estimator applied to a true AR(1) process is flat. A number obtained this way is a property of the analysis window, not of the signal, and no amount of data will make it converge.

Real recordings show both effects, but not in the proportions we expected (\figpanel{fig:windows}{e}). Cutting each recording into non-overlapping windows from 0.5 s upward, and measuring window length in units of the timescale being measured rather than in seconds -- a 1 s window is generous for a 4 ms timescale and hopeless for a 1 s one -- the median estimate turns out to be close to window-independent. In six of the eight recordings it stayed between 0.73 and 1.9 times the estimate from the longest window available, over more than three orders of magnitude of window length. The exceptions are macaque visual ECoG and the mouse spike counts, which sit at about a quarter of their final value until the window exceeds a thousand times the timescale.

This deserves to be stated plainly, because it is not what the same analysis showed before we corrected the autocorrelation. Taking the autocorrelation from the raw signal instead of from the analysis band, estimates collapse by up to two orders of magnitude at the shortest windows and climb steadily with window length, which reads convincingly as a window-length effect on the brain's timescales. It is not. It is drift below the fit band entering the lag domain, and contributing more the longer the window. The window dependence was a property of the estimator. We report it because the corrected version is the weaker and less quotable result, and because an artifact of this shape is easy to produce by accident and hard to notice.

The most consequential result is what happens to the spread between windows (\figpanel{fig:windows}{f}). If finite sampling were the only source of variability, the spread of estimates across windows of a given length would follow the AR(1) sampling bound. It does not. At the longest windows available, the spectral estimate exceeded that bound in all eight recordings, by a median factor of 6 and by up to 30, and the autocorrelation estimate exceeded it in all eight, by a median factor of 7 and by up to 29. Unlike the median, this does not go away when the autocorrelation is corrected. At the durations where a timescale estimate becomes precise enough to be worth reporting, the dominant source of variability is not estimation noise but genuine change in the signal over time, which is the dynamic timescale problem of \figpanel{fig:bias}{c} measured directly. Two practical consequences follow. A confidence interval computed from the AR(1) sampling formula will substantially understate the uncertainty of a windowed estimate, and one computed from the spread across windows will not. And a timescale reported for a long recording is an average over a quantity that was changing throughout it.

\begin{figure}[htbp]
  \centering
  \includesvg[width=\linewidth]{revision/figures/fig_windows.svg}
  \caption{\textbf{How much data a timescale estimate needs.} \textbf{a)} Simulated AR(1) estimates against the ratio of recording length to timescale. Short recordings underestimate, and all three estimators do so similarly. \textbf{b)} The same, for the lagged regression alone, against the analytic mean and its 95\% interval: the finite-sample theory of eq. \ref{eq:finite} predicts the simulated bias without a fitted parameter. \textbf{c)} Spread of the same estimates against the analytic $\sqrt{2\tau/T}$; a 10\% standard error needs about $200\tau$ of data, and precision improves only as the square root of duration. \textbf{d)} The sampling distribution at three recording lengths; the right skew follows from the convexity of eq. \ref{eq:tau}, so a short recording is more likely to underestimate than to overestimate but errs further when it overestimates. \textbf{e)} Real recordings cut into non-overlapping windows, normalized to the estimate from the longest window and plotted against window length in units of that estimate; shading is $\pm 20\%$. Six of the eight stay within a factor of two throughout, macaque ECoG and the mouse spike counts being the exceptions. \textbf{f)} Spread of estimates across windows of a given length, divided by the spread finite sampling alone predicts, at the longest window available. Values above one are variability that sampling cannot explain.}
  \label{fig:windows}
\end{figure}

\section{Discussion}\label{sec:discussion}

A neural timescale is not measured, it is estimated, and it is always the parameter of a model. When that model is a first-order autoregressive process, as it is for an exponential fit to the autocorrelation, a Lorentzian fit to the spectrum, or a lagged regression on the signal itself, then $\tau$ is the single relaxation time of that process and reporting it as a property of a region, a state, or a subject is reasonable. When the model is wrong, $\tau$ is still well defined as a parameter estimate, still precise, still reproducible, and still comparable across conditions. It is simply no longer the quantity the biological interpretation is about. Our simulations show that these two situations look identical from the outside, and our real data show that neural recordings are usually in the second one.

The distinction is checkable, and cheaply. Three diagnostics travel with any timescale estimate at essentially no cost. The first is agreement between independent estimators: when the data are AR(1) the six estimators we compared agree to within 2\%, and when they are not they disagree by up to a factor of 46 in simulation. The second is the fraction of the log spectrum, or of the autocorrelation, left unexplained by the fitted model. The third is the AR order chosen by BIC. If independent estimators agree to within about 1.5-fold, if an AR(1) model leaves only a few percent of the spectrum unexplained, and if BIC selects a low order, then a scalar $\tau$ is defensible. None of the eight recordings we examined met that bar.

We also want to be clear about which failure is which, because they call for different responses. An oscillation added to a signal, such as a stimulus-driven steady-state response, is not an autoregressive process and cannot be modeled by one at any practical order; it must be avoided, and the frequency domain is where avoiding it is easiest, because a narrow peak can be discounted as an outlier while a cosine in the autocorrelation cannot. An oscillation that is part of the dynamics is an AR(2) process, and is best modeled rather than removed; doing so recovers both its rhythm and the decay time of its envelope, which an AR(1) model conflates with the aperiodic timescale. Multiple relaxation processes call for a higher order, and the resulting model is readable as a list of decay times. And a genuinely scale-free process has no timescale to recover, so any number reported for it is a property of the analysis window rather than of the brain.

The finite-sampling result is the most immediately actionable. Because the relative bias and precision of an AR(1) timescale depend on the recording only through the ratio of duration to timescale, the requirement can be stated without reference to sampling rate or to the value of $\tau$: roughly $20\tau$ of data for the bias to fall below 10\%, and roughly $200\tau$ for the standard error to do the same. For a 100 ms cortical timescale that is 2 and 20 seconds of stationary data, per estimate. This is why averaging autocorrelations or spectra across trials, neurons, or channels before fitting works so well \cite{murray_hierarchy_2014, spitzner_mr_2021}: it buys the effective duration that a single epoch cannot. It is also why claims about timescales changing within a trial deserve care, since the window that makes an estimate precise is often longer than the state it is meant to resolve.

Our empirical windowing results add a caution to that arithmetic. In real recordings, the spread of estimates across windows of a given length was roughly an order of magnitude larger than finite sampling alone predicts. The dominant source of variability at usable window lengths is therefore not estimation noise but genuine non-stationarity, exactly the dynamic timescale problem sketched in \figpanel{fig:bias}{c}. A confidence interval computed from the AR(1) sampling formula will substantially understate the uncertainty of a windowed estimate; one computed from the spread across windows will not.

\subsection{Recommendations}

We suggest the following when reporting a neural timescale. \textit{First}, state the model rather than only the estimator, for example ``the relaxation time of an AR(1) fit to the 1-100 Hz spectrum'' rather than ``the neural timescale''. \textit{Second}, report the ratio of recording length to estimated timescale, and treat estimates below roughly $200\tau$ as precision-limited. \textit{Third}, report the fraction of the spectrum or autocorrelation left unexplained by the fitted model. \textit{Fourth}, report at least one independent estimator; disagreement beyond about 1.5-fold is evidence of model mismatch rather than of noise. \textit{Fifth}, if the data contain a strong narrowband rhythm, prefer spectral-domain estimation and state whether the rhythm is being modeled or excluded. \textit{Sixth}, if an AR($p$) model is used, report the modes as $(\tau, f)$ pairs with their importances and the number needed to reproduce the fit, not the coefficient vector. \textit{Finally}, when windowing, take the confidence interval from the spread across windows rather than from the sampling formula.

\subsection{Limitations}

Our simulations treat stationary processes. Non-stationarity is addressed only through the windowing analysis, and it remains the hardest of the four biases to handle, because reducing it by shortening windows increases finite sampling bias. Seven of our eight recordings are continuous field potentials or a haemodynamic signal, and only one is a spike train, for which binning and Poisson variability add a flat floor that an AR spectral model cannot represent and that we handle here only by restricting the fit band \cite{zeraati_flexible_2022, meyer_inferring_2019}; a model with an explicit observation-noise term would treat it properly and we do not attempt one. Our real data analysis deliberately applies one identical pipeline to every recording, so the dataset-specific preprocessing a specialist would apply, such as artifact rejection, burst detection, or oscillation removal, is absent by design; the model mismatch we report is therefore an upper bound on what careful per-dataset treatment would leave. Finally, the mode decomposition is exact for the fitted model but inherits its assumptions. It is a parsimonious description of a spectrum, and whether any given mode corresponds to a separable biological process is a question the model cannot answer. Fitted poles are not latent processes: summing two independent AR(1) processes yields a model whose fitted poles are pulled away from both of the original timescales, because the sum is not itself an autoregressive process.

\subsection{Future}

The most useful next step is probably not another estimator but a shared set of diagnostics. Timescales are now reported across spiking, LFP, ECoG, EEG, MEG and fMRI, on data ranging from seconds to hours, and comparisons across these literatures currently assume that the same word denotes the same quantity. The three checks above are cheap enough to report alongside any estimate, and would make cross-study comparison meaningful in a way that a bare $\tau$ cannot. Beyond that, the mode description suggests a richer target for the hierarchy literature than a single number: a region might be characterized by two or three decay times and the rhythms attached to them, and the question of whether \emph{those} vary systematically along the cortical hierarchy \cite{murray_hierarchy_2014, raut_organization_2020, chaudhuri_large-scale_2015} is one that a single-timescale summary cannot pose. Extending the framework to spike trains, and to models that allow the modes themselves to change over time, would connect it to the dynamic timescale results that motivated much of this work.

\section{Methods}\label{sec:methods}

\subsection{Estimators}\label{sec:methods-est}

Six estimators were compared. They differ in the domain they are fit in and in the model order, and they otherwise share everything: each consumes the same mean-centered signal and the same Welch power spectrum \cite{welch_use_1967}, so that differences between them cannot be attributed to preprocessing.

Making that sharing exact requires care with the autocorrelation, and getting it wrong is easy. A spectral fit is restricted to an analysis band, 1-100 Hz for most of the electrophysiological recordings here, whereas the sample autocorrelation of the raw signal contains everything in the recording, including drift below the band. Fitting one estimator to the band and the other to the whole signal compares two answers to two different questions. We therefore compute the autocorrelation from the same band-limited spectrum the spectral estimators see, by inverse Fourier transform, with bins inside the band that are excluded as line noise filled by log-linear interpolation from their neighbours rather than zeroed, so that notching does not ring in the lag domain. The difference is not cosmetic: on rat hippocampal CA1/CA3 the raw autocorrelation gives an AR(1) timescale of 938 ms, set by slow drift the spectral fit never saw, against 14 ms once both estimators describe the same band.

One further correction is needed. Welch's estimate is the average periodogram of tapered segments, so the autocovariance it implies is the true autocovariance multiplied by the taper's own autocorrelation, which for a Hann window has fallen to 0.66 at a lag one quarter of the segment length. Left in place this taper masquerades as a faster decay and biases every lag-domain estimate downward, so we divide it out, and discard the lags at which the taper retains less than a fifth of its energy, where the correction amplifies noise faster than it removes bias.

In the time domain, the AR(1) estimator is a conditional least-squares regression of $x_t$ on $x_{t-1}$, and the AR($p$) estimator uses Burg's method \cite{burg_maximum_1975}. In the autocorrelation domain, the model is $A\,\rho_p(k) + B$, where $\rho_p$ is the exact AR($p$) autocorrelation, $A$ absorbs a white-noise floor at lag zero and $B$ a constant offset. For $p=1$ this reduces to the exponential-plus-offset model that is standard in this literature \cite{murray_hierarchy_2014}. In the spectral domain, the model is eq. \ref{eq:arp_psd} with a free scale, fit to $\log_{10}$ power, which is the convention used by spectral parameterization methods \cite{donoghue_parameterizing_2020}.

Two choices are shared by the autocorrelation and spectral fits, and both matter for a fair comparison. First, robust (soft-L1) regression is available in both domains and is enabled in both or in neither, never in one alone. It is enabled only when the signal contains a rhythm: in simulation, for the added sinusoid and the AR(2) oscillator and for no other process; in real data, for every recording, since all of them contain a narrowband peak. Robustness is not free. It protects a broadband fit from a narrowband rhythm, which is an outlier in the frequency domain, but when there is no outlier to reject, down-weighting the tails of an ordinary residual distribution biases the fit, so it is a claim about the data and not a default. Second, both are optimized over reflection coefficients rather than over the AR coefficients directly. Every point in the box $|k_i| < 1$ corresponds to a stationary model, the parameters are all of order one, and the fit cannot wander into a region where mode timescales are undefined. Without this reparameterization roughly half of our spectral fits to real data returned non-stationary parameters. Multiple starting points are used because the objective is non-convex when a rhythm is present.

Two further requirements are imposed on the fitted model, and both exist because the model is being used as an instrument for measuring a timescale rather than only as a description of the data. First, the model must contain a relaxation mode. Nothing in a least-squares fit requires this: a converged, well-fitting AR($p$) can place every pole in a conjugate pair, in which case it has no decay time to report at all. This is not rare. Before the constraint, a third of the fits to our real recordings came back with no relaxation mode, having described the data well. We therefore factor one real pole out of the characteristic polynomial, $A(z^{-1}) = (1 - a z^{-1})A_{\mathrm{rest}}(z^{-1})$ with $0 < a < 1$, and parameterize the remainder by reflection coefficients. A relaxation mode then exists by construction, exactly as stationarity does. The requirement costs one order, so it is not applied at order two, where it would leave an AR(2) oscillator unrepresentable. It costs essentially nothing in fit: across our recordings the median change in each estimator's own objective was $+0.003$ for the lag-domain fit and $-0.000$ for the spectral fit.

Spectra are estimated by Welch's method, except for fMRI. A 296-sample record cannot afford Welch's trade: segmenting it leaves either too few frequency bins to fit or too few segments to average. We therefore use a multitaper estimate of the whole record \cite{thomson_spectrum_1982}, with time-bandwidth product $NW = 5$, which gives twice as many bins as the best Welch setting at a quarter of the bin-to-bin variability. Because adjacent multitaper bins inside the resolution bandwidth are not independent, the spectral fits use bins thinned to that spacing; leaving them in makes the record look eight times larger than it is, and order selection then buys orders it cannot support.

Second, the reported aperiodic timescale is the most important \emph{non-rhythmic} mode, where non-rhythmic means $f\tau < 1$, and a real pole qualifies trivially. Two things make this the right choice rather than "the timescale of the real pole". A conjugate pair that rings for a third of a cycle is a damped relaxation, and reporting nothing for a model containing one would be an artifact of the algebra. And the guaranteed real pole is guaranteed to exist, not to matter: on mouse LFP the constrained spectral fit parks its required pole at 0.6 ms with negligible importance while the aperiodic structure sits in an 88 ms pair ringing 0.57 of a cycle. Importance, not pole type, has to decide.

Models fit in either domain are scored in both. Spectral fit is the coefficient of determination of $\log_{10}$ power over the analysis band, with the overall scale refit; autocorrelation fit is the coefficient of determination over the fitted lags, with $A$ and $B$ refit. This matters because the two objectives are not equivalent: a model fit to the autocorrelation of a scale-free process describes that autocorrelation well and its spectrum badly, whereas a model fit to the spectrum describes both.

\subsection{Modes}\label{sec:methods-modes}

The AR($p$) recursion in eq. \ref{eq:arp} has a characteristic polynomial $z^p - \varphi_1 z^{p-1} - \cdots - \varphi_p$, whose roots $\lambda_k$ we call the poles. Stationarity requires $|\lambda_k| < 1$. Dividing by $z^p$ gives
\begin{align}
1 - \sum_{k=1}^{p} \varphi_k z^{-k} = \prod_{k=1}^{p} (1 - \lambda_k z^{-1}), \label{eq:factor}
\end{align}
so the spectral denominator in eq. \ref{eq:arp_psd} factorizes into first-order terms and each pole contributes a decay time and, when complex, a frequency (eq. \ref{eq:mode}). Because the coefficients are real, complex poles come in conjugate pairs, and we combine each pair into a single damped oscillatory mode equivalent to a real AR(2) factor. A real pole is a pure relaxation mode.

Two points about this factorization are easy to state wrongly. It is a factorization of the transfer function, not a decomposition of the signal into independent AR(1) processes: the factors multiply in the spectrum, and an AR($p$) process is a cascade of first-order filters driven by one input. Conversely, the same poles do appear additively in the time domain. The impulse response is $\psi_n = \sum_k c_k \lambda_k^n$, so the response to an input is a sum of decaying modes driven by a shared innovation, and the autocovariance decomposes the same way. Modes therefore add in the autocorrelation and multiply in the spectral denominator, and because they share an input the process variance does not decompose into per-mode shares.

\subsection{Mode Importance}\label{sec:methods-importance}

Writing $g_k(f)$ for the log of the $k$th pole factor, the model log spectrum is $\log \sigma^2 - \sum_k g_k(f)$: modes are additive in log power. Deleting mode $k$ and refitting the overall scale by least squares therefore leaves a residual whose root-mean-square in $\log_{10}$ power is exactly $\mathrm{sd}_f\{g_k\}/\ln 10$, and we use this as the mode's importance. An Itakura-Saito variant, equal to the asymptotic Whittle log-likelihood degradation \cite{whittle_estimation_1953, itakura_analysis_1968}, gives the same ranking in every case we examined. The error of a model retaining an arbitrary subset of modes is a quadratic form in the retained set, so the best subset of each size can be evaluated in closed form; we define the dominant modes as the smallest subset reproducing the fitted log spectrum with $R^2 \ge 0.99$.

\subsection{Model Order}\label{sec:methods-order}

Order is selected by BIC \cite{schwarz_estimating_1978} in the domain the model is fit in: from the conditional Gaussian likelihood in the time domain, and from the residual sum of squares in the autocorrelation and spectral domains. The search runs over a log-spaced grid of orders up to 20 with early stopping, and never fits more parameters than the data support. When the selected order reaches the top of the range this is recorded and reported, because it indicates a process the AR class cannot represent at any practical order rather than a well-determined order.

\subsection{Simulations}\label{sec:methods-sims}

Five processes were used, each with a closed-form spectrum and autocorrelation, simulated at 500 Hz. They are chosen so that each isolates one way the AR(1) assumption can fail, and every one of them appears in the main figure.

\textit{AR(1)} with $\tau=50$ ms is the positive control: the model is correct, so an estimator that fails here fails for reasons of its own. \textit{AR(1) plus sinusoid} adds $A\cos(2\pi f_0 t + \phi)$ with random phase, parameterized by the ratio of oscillatory to aperiodic variance; the sinusoid is not an autoregressive process at any order, and it stands for the rhythm-contaminated recording. \textit{AR(2) oscillator} has poles at $|\lambda| e^{\pm i\theta}$ with $|\lambda| = e^{-1/(\tau f_s)}$ and $\theta = 2\pi f_0/f_s$; it is a rhythm the model represents exactly, and its role is to separate a rhythm that must be avoided from one that should be modeled. \textit{AR($p$)} is a constructed AR(8) with five modes, a 100 ms and a 2 ms relaxation and damped rhythms at 10, 40 and 180 Hz. It is the only process here that is inside the model family but above first order, which is what makes it necessary: every other failing process confounds model mismatch with order mismatch, and this one isolates the second. It is the same model read mode by mode in \figpanel{fig:modes}{}, so the benchmark and that figure describe one process. \textit{Power law} is a long-memory ARFIMA process \cite{granger_introduction_1980, hosking_fractional_1981} simulated exactly by circulant embedding \cite{davies_tests_1987}; it is outside the family and has no characteristic timescale at all, so there is nothing for an estimator to recover.

For the AR(8) the target is the true model's own aperiodic timescale, obtained by applying to the true poles the same rule an estimator applies to fitted ones -- the most important non-rhythmic mode. Defining it instead as the slowest relaxation by construction would score estimators against a quantity they were never asked for, since the 10 Hz mode of that model rings only half a cycle and is both non-rhythmic by that rule and more important than the 100 ms relaxation. Every theoretical spectrum and autocorrelation was verified against long simulations.

\subsection{Finite Sampling}\label{sec:methods-finite}

For a mean-corrected least-squares AR(1) fit on $n$ samples, the classical results are $\mathbb{E}[\hat\varphi] - \varphi \approx -(1+3\varphi)/n$ \cite{marriott_bias_1954, kendall_note_1954} and eq. \ref{eq:var}. Propagating both through eq. \ref{eq:tau} to second order, and writing $T = n/f_s$ for the recording duration, gives in the regime $\tau f_s \gg 1$ that all neural recordings occupy
\begin{align}
\frac{\mathbb{E}[\hat\tau] - \tau}{\tau} \approx -\frac{2\tau}{T}, \qquad
\frac{\mathrm{sd}[\hat\tau]}{\tau} \approx \sqrt{\frac{2\tau}{T}}. \label{eq:finite}
\end{align}
Both depend on the recording only through $T/\tau$. Confidence intervals are obtained by mapping the interval for $\hat\varphi$ through eq. \ref{eq:tau}, which preserves coverage because the map is monotone; refinements for $\varphi$ very close to one are available \cite{gallagher_small_2008}.

\subsection{Datasets}\label{sec:methods-data}

Eight publicly available recordings were used, selected to span species, recording modality and oscillatory regime rather than to address any single biological question. \textit{fMRI}: ABIDE resting-state BOLD \cite{di_martino_autism_2014}, six participants from the UM site, preprocessed by C-PAC without temporal filtering, TR 2 s. We take the precuneus, left and right, giving twelve traces. Restricting to one named region matters here in a way it does not for the electrophysiological recordings. A single 296-sample BOLD run gives a spectrum too noisy to fit, and averaging over regions is the obvious remedy but the wrong one: regions differ in exactly the property being measured, and precuneus, medial prefrontal and occipital spectra have visibly different shapes in these data, so an average over all 111 Harvard-Oxford regions would describe no region at all. One default-mode hub keeps the average interpretable and makes a trace here the same kind of object as a channel elsewhere. The unfiltered pipeline is used deliberately, since the usual 0.01-0.1 Hz band-pass removes most of the spectrum a timescale is estimated from. BOLD is the hardest case for spectral estimation in this collection, and it is included for that reason: a 592 s run sampled every 2 s contains 296 points, so Welch's trade of resolution against variance has nothing to spend. Segmenting leaves either too few frequency bins to fit or too few segments to average, which is why we use a multitaper estimate of the whole record instead (Methods \ref{sec:methods-est}). \textit{MEG}: the FieldTrip cortico-muscular coherence recording \cite{schoffelen_neuronal_2005, oostenveld_fieldtrip_2011}, eight sensorimotor magnetometers during isometric contraction. \textit{EEG}: MPI-LEMON eyes-closed rest \cite{babayan_mind_2019}, posterior channels. \textit{ECoG}: pre-DBS primary motor cortex in Parkinson disease \cite{cole_nonsinusoidal_2017}; and macaque visual cortex during a grating task \cite{nagasaka_multidimensional_2011}. \textit{LFP}: mouse Neuropixels in VISam during natural movie viewing \cite{siegle_survey_2021}; and rat hippocampal CA1/CA3 \cite{mizuseki_theta_2009, mizuseki_hc3_2013}. \textit{Spiking}: spike counts of eight VISp units from the same Neuropixels session \cite{siegle_survey_2021}, binned at 50 ms over the same 300 s of natural movie. These are fitted over 0.05-2 Hz rather than the 0.05-8 Hz the sampling rate allows, because above about 1.5 Hz the binned-count spectrum is within a factor of 1.5 of its flat Poisson shot-noise floor, which $\sigma^2/|A|^2$ has no term to represent; the criterion is that the process remain above its own shot noise, not a fixed frequency. Every other recording here is a continuous field potential or a haemodynamic signal, while the neural-timescale literature rests largely on spike-count autocorrelations, so this is the one modality whose absence would be conspicuous; it also contributes a second mouse visual region.

Electrophysiological recordings were resampled to a common 250 Hz and fit over 1-100 Hz, with line-noise fundamentals and harmonics excluded from the spectral fits. fMRI carries its own analysis band, 0.01-0.24 Hz, because it is sampled three orders of magnitude more slowly. No other filtering, artifact rejection or oscillation removal was applied. For the windowing analysis, each recording was cut into non-overlapping windows of increasing length and every window analysed with the identical pipeline.

\subsection{Code}\label{sec:methods-code}

All analyses use an open-source Python module with 133 unit tests covering, among other things, exact reconstruction of the AR($p$) spectrum from its pole factorization, agreement between the closed-form mode importance and a brute-force deletion computation, agreement of every theoretical spectrum and autocorrelation with long simulations, and agreement of eq. \ref{eq:finite} with Monte Carlo. Code, figure scripts and result tables are available at \href{https://github.com/voytekresearch/timescale-methods}{https://github.com/voytekresearch/timescale-methods}.

\backmatter

\clearpage
\section*{Supplementary Information}\label{sec:supp}

The main text emphasizes what the results mean for reporting a neural timescale. Conventions, derivations, and the full parameter sweeps are collected below.

\setcounter{figure}{0}
\renewcommand{\thefigure}{S\arabic{figure}}
\setcounter{table}{0}
\renewcommand{\thetable}{S\arabic{table}}

\subsection*{Conventions}

We write the AR($p$) model as in eq. \ref{eq:arp}, so that the characteristic polynomial is $z^p - \varphi_1 z^{p-1} - \cdots - \varphi_p$ and stationarity requires all of its roots to lie inside the unit circle. Sign conventions differ between texts and the two common polynomials are reciprocals of one another; we use this one throughout, so that ``pole inside the unit circle'' means ``stationary''. The one-sided spectral density is
\begin{align}
S(f) = \frac{2\sigma^2/f_s}{\left|1 - \sum_{k=1}^{p}\varphi_k e^{-i 2\pi f k/f_s}\right|^{2}}, \qquad 0 \le f \le f_s/2,
\end{align}
normalized so that its integral is the process variance, matching Welch's method \cite{welch_use_1967}. For $p=1$ this is the discrete-time Lorentzian whose knee frequency is $1/(2\pi\tau)$. The theoretical autocovariance is obtained by solving the Yule-Walker system for the first $p+1$ lags and recursing thereafter.

\subsection*{The Integrated Timescale}

As a model-free reference we use $\tau_{\mathrm{int}} = \int_0^\infty \rho(t)\,dt$ \cite{sokal_monte_1997, priestley_spectral_1981}, evaluated with the trapezoid convention so that a sampled exponential returns its continuous-time $\tau$ rather than $\tau + 1/(2f_s)$. Equivalently $\tau_{\mathrm{int}} = S(0)/(4\gamma(0))$. It equals $\tau$ for an AR(1) process, is the variance-weighted mean of component timescales for a sum of independent AR(1) processes, and diverges for a long-memory process. It is not the envelope decay time of an oscillator: for a damped oscillator the anticorrelated troughs of the autocorrelation cancel most of the integral, which is why the AR(2) simulation in the main text has $\tau_{\mathrm{int}} = 9$ ms and an envelope decay of 50 ms.

\subsection*{Stationarity by Construction}

Fitting AR coefficients directly requires bounds that grow like binomial coefficients, and a bounded optimizer routinely returns non-stationary parameters for which mode timescales are undefined. We instead optimize over reflection coefficients $k_i$, related to the AR coefficients by the Levinson recursion. The condition $|k_i| < 1$ for all $i$ is equivalent to stationarity, so the entire search box is admissible and all parameters are of order one. In our real data, 48\% of spectral fits returned non-stationary parameters when fit over the coefficients directly, and none did after reparameterization.

An alternative repair, useful when a fit is inherited from elsewhere, is to reflect any pole with $|\lambda| \ge 1$ to $1/\bar\lambda$ and divide the scale by $|\lambda|^2$. On the unit circle $|1 - \lambda z^{-1}|^2 = |\lambda|^2 |1 - z^{-1}/\bar\lambda|^2$, so this leaves the model spectrum exactly unchanged while making the process stationary.

\subsection*{Supplementary Figures}

\begin{figure}[htbp]
  \centering
  \includesvg[width=\linewidth]{revision/figures/figS_oscillation_grid.svg}
  \caption{\textbf{Full oscillation sweep.} Median estimated aperiodic timescale divided by the truth, across oscillation frequency and oscillation power, for the four fitted estimators. One is exact and both the printed number and the colour say so; the colour scale is spaced logarithmically, so a factor of two too large and a factor of two too small are equally far from white. Every cell is filled, because the fit is constrained to contain a relaxation mode (Methods \ref{sec:methods-est}); before that constraint the lag-domain fits left cells empty, having spent every pole on the rhythm. Both spectral fits stay between 0.97 and 1.03 everywhere, which is why those two panels are uniformly white: robust regression discards the rhythm and leaves the aperiodic fit essentially untouched.}
  \label{fig:s_grid}
\end{figure}

\bibliography{references,references_revision}

\end{document}
